\vspace*{\fill}
\begin{glyph}{Capital (京)}{capital}\label{glyph:capital}
Capital; Large metropolis.
\end{glyph}

\begin{comps}
\comprtwocone & 亠 + 口 + 小  \\\hline
\comprthreecone & Head/Above + Mouth (\ref{glyph:mouth}) + Small (\ref{glyph:small}) \\\hline
\end{comps}

\begin{remglyph}
\remcom{Small} \remcom{vampire} \remcom{heads} in the \hooglig{capital}.\\\hline
\end{remglyph}

\begin{prons}
\pronsrtwocone & \textbf{キョウ (KY\={O})} & ---\\\hline
\pronsrthreecone & ケイ (KEI) & --- \\\hline
\end{prons}

\begin{remprons}
The \comon{Kyoto} \hooglig{capital} \ucomon{caters} for all your needs.\\\hline
\end{remprons}

\begin{somme}
上京 & upper + capital = \textit{come/go to Tokyo} & ジョウ{\middel}キョウ (J\={O}{\middel}KY\={O})\\\hline
東京 & east + capital = \textit{Tokyo} & トウ{\middel}キョウ (T\={O}{\middel}KY\={O})\\\hline
\end{somme}

\begin{sentence}
\ruby{\underline{中山}}{なかやま}\underline{さん}\ruby{は}{わ}\ruby{\underline{上京}}{ジョウキョウ}しました。\\\hline
Naka{\middel}yama-san wa J\={O}{\middel}KY\={O} shimashita.\\\hline
(Nakayama-san) (went to Tokyo)\\\hline
=\textit{Nakayama-san went to Tokyo.}\\\hline
\end{sentence}
\end{multicols}
\vspace*{-\baselineskip}
\vhrulefill{5pt}
%%--------------------------------------------------------------------
%%--------------------------------------------------------------------
\begin{multicols}{2}
\begin{glyph}{See (見)}{see}\label{glyph:see}
See; Watch.\\\hline
The idea of ``show'' can also be expressed.
\end{glyph}

\begin{comps}
\comprtwocone & 目 + 儿  \\\hline
    \comprthreecone & Eye (\ref{glyph:eye}) + man/a person (at the bottom of a
    character)\\\hline
\end{comps}

\begin{remglyph}
    \remcom{People} with \remcom{eyes} can \hooglig{see}.\\\hline
\end{remglyph}

\begin{prons}
\pronsrtwocone & \textbf{ケン (KEN)} & \textbf{み (mi)}\\\hline
\pronsrthreecone & --- & --- \\\hline
\end{prons}

\begin{remprons}
I \hooglig{see} the \comon{kennels} have \comkun{mirrors}.\\\hline
\end{remprons}

\begin{somme}
見える (intr) & \textit(to be visible) & み{\middel}える (mi{\middel}eru)\\\hline
見る (tr) & \textit{to see; to watch)} & み{\middel}る (mi{\middel}ru)\\\hline
見せる (tr) & \textit{to show (something)} & み{\middel}せる (mi{\middel}seru)\\\hline
見上げる & see + upper = \textit{to look up (to/at)} & み　あ{\middel}げる (mi a{\middel}geru)\\\hline
月見 & moon + see = \textit{moon viewing} & つき{\middel}み (tsuki{\middel}mi)\\\hline
見失う & see + lose = \textit{to lose sight of} & み{\middel}うしな{\middel}う (mi{\middel}ushina{\middel}u)\\\hline
外見 & outside + see = \textit{external appearance} & ガイ{\middel}ケン (GAI{\middel}KEN)\\\hline
\end{somme}

\begin{sentence}
\ruby{\underline{山}}{やま}の\ruby{\underline{上}}{うえ}に \ruby{\underline{小}}{ちい}\underline{さい}{\;}\ruby{\underline{木}}{き}が\ruby{\underline{見}}{み}\underline{えます}。\\\hline
yama no ue ni chii{\middel}sai ki ga mi{\middel}emasu.\\\hline
(mountain) (upper) (small) (tree) (is visible)\\\hline
=\textit{A small tree is visible on top of the mountain.}\\\hline
\end{sentence}
\end{multicols}
\vspace*{-\baselineskip}
\vhrulefill{5pt}
\vspace*{\fill}
%%--------------------------------------------------------------------
%%--------------------------------------------------------------------
\pagebreak
\vspace*{\fill}
\begin{multicols}{2}
\begin{glyph}{Stop (止)}{stop}\label{glyph:stop}
Stop.
\end{glyph}

\begin{comps}
\comprtwocone & ⺊ + 丨 + 一\\\hline
\comprthreecone & Divination + Pole + One (\ref{glyph:one}) \\\hline
\end{comps}

\begin{remglyph}
    \hooglig{Stop} using that \remcom{divination pole} at \remcom{once}.\\\hline
\end{remglyph}

\begin{prons}
\pronsrtwocone & \textbf{シ (SHI)} & \textbf{と (to)}\\\hline
\pronsrthreecone & --- & --- \\\hline
\end{prons}

\begin{remprons}
\hooglig{Stop} the \comkun{torrent} with a \comon{shield}.\\\hline
\end{remprons}

\begin{somme}
止まる (intr) & \textit{to stop} & と{\middel}まる (to{\middel}maru)\\\hline
止める (tr) & \textit{to stop (an object)} & と{\middel}める (to{\middel}meru)\\\hline
中止 & middle + stop = \textit{discontinue} & チュウ{\middel}シ (CH\={U}{\middel}SHI)\\\hline
立ち止まる & stand + stop = \textit{discontinue} & た{\middel}ち　ど{\middel}まる (ta{\middel}chi do{\middel}maru)\\\hline
車止め & car + stop = \textit{closed to vehicles} & くるま　ど{\middel}め (kuruma do{\middel}me)\\\hline
\end{somme}

\begin{sentence}
\underline{こちら}に\ruby{\underline{止}}{と}\underline{まって} \ruby{\underline{下}}{くだ}\underline{さい}。\\\hline
kochira ni to{\middel}matte kuda{\middel}sai.\\\hline
(here) (stop) (please)\\\hline
=\textit{Please stop here.}\\\hline
\end{sentence}
\end{multicols}
\vspace*{-\baselineskip}
\vhrulefill{5pt}
%%--------------------------------------------------------------------
%%--------------------------------------------------------------------
\begin{multicols}{2}
\begin{glyph}{Seven (七)}{seven}\label{glyph:seven}
Seven.
\end{glyph}

\begin{comps}
\comprtwocone & 一 + 亅\\\hline
\comprthreecone & One + Hook \\\hline
\end{comps}

\begin{remglyph}
    \remcom{Captain Hook} sailed the \hooglig{seven} seas \remcom{once} in his entire life.\\\hline
\end{remglyph}

\begin{prons}
\pronsrtwocone & \textbf{シチ (SHICHI)} & \textbf{なな (nana)}\\\hline
\pronsrthreecone & --- & なの (nano) \\\hline
\end{prons}

\begin{remprons}
\comon{She cheated} the \comkun{nana} with \hooglig{seven} \ucomkun{nanoparticles}.\\\hline
\end{remprons}

\begin{somme}
七 & \textit{seven} & シチ (SHICHI)\\\hline
七つ & \textit{seven (general counter)} & なな{\middel}つ (nana{\middel}tsu)\\\hline
七月 & seven + moon (month) \textit{July} & シチ{\middel}ガツ (SHICHI{\middel}GATSU)\\\hline
十七 & ten + seven = \textit{seventeen} & ジュウ{\middel}なな (J\={U}{\middel}nana)\\\hline
七十 & seven + ten = \textit{seventy} & なな{\middel}ジュウ (nana{\middel}J\={U})\\\hline
七時 & seven + time = \textit{seven o'clock} & シチ{\middel}ジ (SHICHI{\middel}JI)\\\hline
\end{somme}

\begin{sentence}
\ruby{\underline{高木}}{たかぎ}\underline{さん}が\ruby{\underline{七月}}{シチガツ}に\ruby{\underline{中米}}{チュウベイ}から\ruby{\underline{来}}{き}\underline{ます}。\\\hline
Taka{\middel}gi-san ga SHICHI{\middel}GATSU ni CH\={U}{\middel}BEI kara ki{\middel}masu.\\\hline
(Takagi-san) (July) (Central America) (will come)\\\hline
=\textit{Takagi-san will come from Central America in July.}\\\hline
\end{sentence}
\end{multicols}
\vspace*{-\baselineskip}
\vhrulefill{5pt}
\vspace*{\fill}
%%--------------------------------------------------------------------
%%--------------------------------------------------------------------
\pagebreak
\vspace*{\fill}
\begin{multicols}{2}
\begin{glyph}{Few (少)}{few}\label{glyph:few}
Few; Little (in terms of quantity).
\end{glyph}

\begin{comps}
\comprtwocone & 小 + ノ \\\hline
    \comprthreecone & Small (\ref{glyph:small}) + NO (katakana)\\\hline
\end{comps}

\begin{remglyph}
A \hooglig{few} occurences of \remcom{small} \remcom{noses}.\\\hline
\end{remglyph}

\begin{prons}
\pronsrtwocone & \textbf{ショウ (SH\={O})} & \textbf{すく、すこ (suku, suko)}\\\hline
\pronsrthreecone & --- & --- \\\hline
\rye{3}{ショウ will appear in compounds.\\\hline
すく and すこ are only found in the first and second examples below.}
\end{prons}

\begin{remprons}
He brought a \hooglig{few} \comkun{scones} with his \comkun{scooter} to the \comon{shore}.\\\hline
\end{remprons}

\begin{somme}
少ない & \textit{few} & すく{\middel}ない (suku{\middel}nai)\\\hline
少し & \textit{a little; a few} & すこ{\middel}し (suko{\middel}shi)\\\hline
少々 & few + few = \textit{a little} & ショウ{\middel}ショウ (SH\={O}{\middel}SH\={O})\\\hline
少女 & few + woman = \textit{girl} & ショウ{\middel}ジョ (SH\={O}{\middel}JO)\\\hline
少年 & few + year = \textit{boy} & ショウ{\middel}ネン (SH\={O}{\middel}NEN)\\\hline
最小 & most + few = \textit{fewest} & サイ{\middel}ショウ (SAI{\middel}SH\={O})\\\hline
\end{somme}

\begin{sentence}
\underline{この} \ruby{\underline{家}}{いえ}に\ruby{は}{わ}\ruby{\underline{物}}{もの}が\ruby{\underline{少}}{すく}\underline{ない}。\\\hline
kono ie ni wa mono ga suku{\middel}nai.\\\hline
(this) (house) (thing) (few)\\\hline
=\textit{There are few things in this house.}\\\hline
\end{sentence}
\end{multicols}
\vspace*{-\baselineskip}
\vhrulefill{5pt}
%%--------------------------------------------------------------------
%%--------------------------------------------------------------------
\begin{multicols}{2}
\begin{glyph}{South (南)}{south}\label{glyph:south}
South.
\end{glyph}

\begin{comps}
\comprtwocone & 十 + 冂 + 䒑 + 十\\\hline
\comprthreecone & Ten (\ref{glyph:ten}) + Wilderness + Grass + Ten (\ref{glyph:ten})\\\hline
\end{comps}

\begin{remglyph}
    The \hooglig{south} \remcom{wilderness} do not even have \remcom{ten}
    \remcom{grass} blades.\\\hline
\end{remglyph}

\begin{prons}
\pronsrtwocone & \textbf{ナン (NAN)} & \textbf{みなみ (minami)}\\\hline
\pronsrthreecone & --- & --- \\\hline
\end{prons}

\begin{remprons}
The \hooglig{southern} \comon{nun} at a \comkun{mean of a meal}.\\\hline
\end{remprons}

\begin{somme}
南 & \textit{south} & みなみ (minami)\\\hline
南口 & south + mouth = \textit{southern exit/entrance} & みなみ{\middel}ぐち (minami{\middel}guchi)\\\hline
南米 & south + rice = \textit{Southern America} & ナン{\middel}ベイ (NAN{\middel}BEI)\\\hline
最南 & most + south = \textit{southernmost} & サイ{\middel}ナン (SAI{\middel}NAN)\\\hline
南西 & south + west = \textit{southwest} & ナン{\middel}セイ (NAN{\middel}SEI)\\\hline
南東 & south + east = \textit{south east} & ナン{\middel}トウ (NAN{\middel}T\={O})\\\hline
\end{somme}

\begin{sentence}
\underline{バスターミナル}の\ruby{\underline{南口}}{みなみぐし}で\ruby{\underline{会}}{あ}\underline{いましょう}。\\\hline
basut\={a}minaru no minami{\middel}guchi de a{\middel}imash\={o}.\\\hline
(bus terminal) (south entrance) (let's meet)\\\hline
=\textit{Let's meet at the south entrance of the bus terminal.}\\\hline
\end{sentence}
\end{multicols}
\vspace*{-\baselineskip}
\vhrulefill{5pt}
\vspace*{\fill}
%%--------------------------------------------------------------------
%%--------------------------------------------------------------------
\pagebreak
\vspace*{\fill}
\begin{multicols}{2}
\begin{glyph}{Craft (工)}{craft}\label{glyph:craft}
Carft; Construction; Manufacturing; Handicrafts; Industrial products.\\\hline
When found at the end of a compound (as in the final example below), this kanji can take on the meaning of ``worker'', and usually implies a person involved in some part of manual labor.
\end{glyph}

\begin{comps}
\comprtwocone & 工\\\hline
\comprthreecone & Craft (\ref{glyph:craft})\\\hline
\end{comps}

\begin{remglyph}
Anvil.\\\hline
\end{remglyph}

\begin{prons}
\pronsrtwocone & \textbf{コウ (K\={O})} & ---\\\hline
\pronsrthreecone & ク (KU) & --- \\\hline
\end{prons}

\begin{remprons}
His \ucomon{cooking} brilliance \comon{caused} me to refine my own \hooglig{craft}.\\\hline
\end{remprons}

\begin{somme}
人工 & person + craft = \textit{artificial}  & ジン{\middel}コウ (JIN{\middel}K\={O})\\\hline
人工林 & person + craft + grove = \textit{planted forest} & ジン{\middel}コウ{\middel}リン (JIN{\middel}K\={O}{\middel}RIN)\\\hline
工学 & craft + study = \textit{engineering} & コウ{\middel}ガク (K\={O}{\middel}GAKU)\\\hline
工学士 & craft + study + gentleman = \textit{Bachelor of Engineering} & コウ{\middel}ガク{\middel}シ (K\={O}{\middel}GAKU{\middel}SHI)\\\hline
刀工 & sword + craft = \textit{swordmaker} & トウ{\middel}コウ (T\={O}{\middel}K\={O})\\\hline
\end{somme}

\begin{sentence}
\underline{これ}\ruby{は}{わ}\ruby{\underline{人工}}{ジンコウ}の\ruby{\underline{島}}{しま}{\;}\underline{ですね}。\\\hline
kore wa JIN{\middel}K\={O} no shima desu ne.\\\hline
(this) (artificial) (island) (isn't it)\\\hline
=\textit{This is an artificial island, isn't it?}\\\hline
\end{sentence}
\end{multicols}
\vspace*{-\baselineskip}
\vhrulefill{5pt}
%%--------------------------------------------------------------------
%%--------------------------------------------------------------------
\begin{multicols}{2}
\begin{glyph}{Left (左)}{hidari}\label{glyph:left}
Left; Left-hand side.
\end{glyph}

\begin{comps}
\comprtwocone & 一 + ノ + 工\\\hline
\comprthreecone & One (\ref{glyph:one}) + NO ノ Katakana + Craft (\ref{glyph:craft})\\\hline
\end{comps}

\begin{remglyph}
The \remcom{superhero} \hooglig{left} the \remcom{anvil} on the monster's head.\\\hline
\end{remglyph}

\begin{prons}
\pronsrtwocone & --- & \textbf{ひだり (hidari)}\\\hline
\pronsrthreecone & サ (SA) & --- \\\hline
\end{prons}

\begin{remprons}
We \hooglig{left} the \comon{subject} of where the \comkun{ring was hid}.\\\hline
\end{remprons}

\begin{somme}
左 & \textit{left} & ひだり (hidari)\\\hline
左回り & left + rotate = \textit{counter-clockwise} & ひだり　まわ{\middel}り (hidari mawa{\middel}ri)\\\hline
左上 & left + upper = \textit{upper left} & ひだり{\middel}うえ (hidari{\middel}ue)\\\hline
左下 & left + lower = \textit{lower left} & ひだり{\middel}した (hidari{\middel}shita)\\\hline
左手 & left + hand = \textit{left hand} & ひだり{\middel}て (hidari{\middel}te)\\\hline
左足 & left + leg = \textit{left leg} & ひだり{\middel}あし (hidari{\middel}ashi)\\\hline
\end{somme}

\begin{sentence}
\underline{あの} \ruby{\underline{人}}{ひと}の\ruby{\underline{左手}}{ひだりて}\ruby{を}{お}\ruby{\underline{見}}{み}\underline{て} \ruby{\underline{下}}{くだ}\underline{さい}。\\\hline
ano hito no hidari{\middel}te o mi{\middel}te kuda{\middel}sai.\\\hline
(that) (person) (left hand) (see) (please)\\\hline
=\textit{Please watch that person's left hand.}\\\hline
\end{sentence}
\end{multicols}
\vspace*{-\baselineskip}
\vhrulefill{5pt}
\vspace*{\fill}
%%--------------------------------------------------------------------
%%--------------------------------------------------------------------
\pagebreak
\begin{multicols}{2}
\begin{glyph}{Tall (高)}{tall}\label{glyph:tall}
Tall; High; Best.\\\hline
This kanji is also used to express the idea of expensive.
\end{glyph}

\begin{comps}
\comprtwocone & 亠 + 口 + 冂 + 口 \\\hline
\comprthreecone & above/head + Mouth (\ref{glyph:mouth}) + wilderness + Mouth(\ref{glyph:mouth})\\\hline
\end{comps}

\begin{remglyph}
    The \hooglig{tall} \remcom{vampire's} \remcom{head} was cast into the
    \remcom{wilderness}.\\\hline
\end{remglyph}

\begin{prons}
\pronsrtwocone & \textbf{コウ (K\={O})} & \textbf{たか (taka)}\\\hline
\pronsrthreecone & --- & --- \\\hline
\end{prons}

\begin{remprons}
I \comkun{tucked} the \hooglig{tall} stash under the bed \comon{'cause} I was afraid.\\\hline
\end{remprons}

\begin{somme}
高い & \textit{tall; expensive} & たか{\middel}い (taka{\middel}i)\\\hline
高まる (intr) & \textit{to rise} & たか{\middel}まる (taka{\middel}maru)\\\hline
高める (tr) & \textit{to raise} & たか{\middel}める (taka{\middel}meru)\\\hline
高山 & tall + mountain = \textit{alpine} & コウ{\middel}ザン (K\={O}{\middel}ZAN)\\\hline
高校 & tall + school = \textit{high school} & コウ{\middel}コウ (K\={O}{\middel}K\={O})\\\hline
円高 & circle (yen) + tall = \textit{strong yen} & エン{\middel}だか (EN{\middel}daka)\\\hline
最高 & most + tall = \textit{highest; the best} & サイ{\middel}コウ (SAI{\middel}K\={O})\\\hline
\end{somme}

\begin{sentence}
\underline{あの} \ruby{\underline{高}}{たか}\underline{い} \ruby{\underline{山}}{やま}\ruby{は}{わ}\ruby{\underline{美}}{うつく}\underline{しい}{\;}\underline{ですね}。\\\hline
ano takai yama wa utsuku{\middel}shii desu ne.\\\hline
(that) (tall) (mountain) (beautiful) (isn't it)\\\hline
=\textit{That tall mountain is beautiful, isn't it.}\\\hline
\end{sentence}
\end{multicols}
\vspace*{-\baselineskip}
\vhrulefill{5pt}
%%--------------------------------------------------------------------
%%--------------------------------------------------------------------
\begin{multicols}{2}
\begin{glyph}{Buy (買)}{buy}\label{glyph:buy}
Buy.
\end{glyph}

\begin{comps}
\comprtwocone & 罒 + 貝 \\\hline
\comprthreecone & net + Shellfish (\ref{glyph:shellfish})\\\hline
\end{comps}

\begin{remglyph}
\hooglig{Buy} the \remcom{shellfish} freshly caught in the \remcom{net}.\\\hline
\end{remglyph}

\begin{prons}
\pronsrtwocone & --- & \textbf{か (ka)}\\\hline
\pronsrthreecone & バイ (BAI) & --- \\\hline
\end{prons}

\begin{remprons}
A \comkun{customary} \hooglig{\ucomon{buy}}.\\\hline
\end{remprons}

\begin{somme}
買う & \textit{to buy} & か{\middel}う (ka{\middel}u)\\\hline
買い上げる & buy + upper = \textit{to buy (up/out)} & か{\middel}い　あ{\middel}げる (ka{\middel}i a{\middel}geru)\\\hline
買上げ & buy + upper = \textit{a purchase} & かい　あ{\middel}げ (kai a{\middel}ge)\\\hline
買い物 & buy + thing = \textit{shopping} & か{\middel}い　もの (ka{\middel}i mono)\\\hline
買い入れる & buy + enter = \textit{to stock up on} & か{\middel}い　い{\middel}れる (ka{\middel}i i{\middel}reru)\\\hline
買い手 & buy + hand = \textit{buyer} & か{\middel}い　て (ka{\middel}i te)\\\hline
\rye{3}{By comparing the second and third examples above, we se that the noun in example 3 has ``lost'' the い that would be present if the kanji were in its -ます verb form, かいます (買います).  It is in a sense implied, based on the hiragana ending after 上.  Unfortunately, there are no hard and fast rules for this aspect of the language, and it can be frustrating when such words are encountered.  Example 3 can be seen variously as 買い上げ、買上げ、or even 買上 in different dictionaries.\\\hline
Nonetheless, not many kanji present with such difficulties, and once a `primary' word (usually the -ます form of the verb) has been learned, its related nounds will be easily recognized.}
\end{somme}

\begin{sentence}
\ruby{\underline{山中}}{やまなか}\underline{さん}\ruby{は}{わ}\ruby{\underline{高}}{たか}\underline{い} \ruby{\underline{物}}{もの}\ruby{を}{お}\ruby{\underline{買}}{か}\underline{う} \underline{つもり}です。\\\hline
Yama{\middel}naka-san wa taka{\middel}i mono o ka{\middel}u tsumori desu.\\\hline
(Yamanaka-san) (expensive) (thing) (buy) (plans)\\\hline
=\textit{Yamanaka-san plans to buy something expensive.}\\\hline
\end{sentence}
\end{multicols}
\vspace*{-\baselineskip}
\vhrulefill{5pt}
%%--------------------------------------------------------------------
%%--------------------------------------------------------------------
\pagebreak
\vspace*{\fill}
\begin{multicols}{2}
\begin{glyph}{Hundred (百)}{hundred}\label{glyph:hundred}
Hundred.\\\hline
In a few instances, the idea of `many' is conveyed.
\end{glyph}

\begin{comps}
\comprtwocone & 一 + 白 \\\hline
\comprthreecone & One (\ref{glyph:one}) + White (\ref{glyph:white})\\\hline
\end{comps}

\begin{remglyph}
\hooglig{Hundred} percent of \remcom{top buns} are \remcom{white}.\\\hline
\end{remglyph}

\begin{prons}
\pronsrtwocone & \textbf{ヒャク (HYAKU)} & ---\\\hline
\pronsrthreecone & --- & --- \\\hline
\end{prons}

\begin{remprons}
\hooglig{Hundred} percent \comon{yucky}.\\\hline
\end{remprons}

\begin{somme}
百 & \textit{one hundred} & ヒャク (HYAKU)\\\hline
二百 & two + hundred = \textit{two hundred} & ニ{\middel}ヒャク (NI{\middel}HYAKU)\\\hline
三百 & three + hundred = \textit{three hundred} & サン{\middel}ビャク (SAN{\middel}BYAKU)\\\hline
四百 & four + hundred = \textit{four hundred} & よん{\middel}ヒャク (yon{\middel}HYAKU)\\\hline
五百 & five + hundred = \textit{five hundred} & ゴ{\middel}ヒャク (GO{\middel}HYAKU)\\\hline
六百 & six + hundred = \textit{six hundred} & ロク{\middel}ヒャク (ROKU{\middel}HYAKU)\\\hline
七百 & seven + hundred = \textit{seven hundred} & なな{\middel}ヒャク (nana{\middel}HYAKU)\\\hline
八百 & eight + hundred = \textit{eight hundred} & ハチ{\middel}ヒャク (HACHI{\middel}HYAKU)\\\hline
九百 & nine + hundred = \textit{nine hundred} & キュウ{\middel}ヒャク (KY\={U}{\middel}HYAKU)\\\hline
\end{somme}

\begin{sentence}
\ruby{\underline{肉}}{ニク}\ruby{を}{お}\ruby{\underline{五百}}{ゴヒャク}{\;}\underline{グラム}{\;}\ruby{\underline{買}}{か}\underline{いました}。\\\hline
NIKU o GO{\middel}HYAKU guramu ka{\middel}imashita.\\\hline
(meat) (five hundred) (grams) (bought)\\\hline
=\textit{(I) bought five hundred grams of meat.}\\\hline
\end{sentence}
\end{multicols}
\vspace*{-\baselineskip}
\vhrulefill{5pt}
%%--------------------------------------------------------------------
%%--------------------------------------------------------------------
\begin{multicols}{2}
\begin{glyph}{Circle (円)}{circle}\label{glyph:circle}
Circle; Round.\\\hline
This kanji is also used to denote the Yen, Japan's Currency.  The connection originated from round one-yen coins.
\end{glyph}

\begin{comps}
\comprtwocone & 冂 + 亠 \\\hline
\comprthreecone & Wilderness + above/head\\\hline
\end{comps}

\begin{remglyph}
    \remcom{Above} the \remcom{wilderness} is a \hooglig{circular} sun.\\\hline
\end{remglyph}

\begin{prons}
\pronsrtwocone & \textbf{エン (EN)} & \textbf{まる (maru)}\\\hline
\pronsrthreecone & --- & --- \\\hline
\end{prons}

\begin{remprons}
The \hooglig{round} \comkun{Maru} cat made quite an \comon{entrance}.\\\hline
\end{remprons}

\begin{somme}
円い & \textit{circular} & まる{\middel}い (maru{\middel}i)\\\hline
円高 & yen + tall = \textit{strong yen} & エン{\middel}だか (EN{\middel}daka)\\\hline
二円 & two + yen = \textit{two yen} & ニ{\middel}エン (NI{\middel}EN)\\\hline
半円 & haf + circle = \textit{semi-circle} & ハン{\middel}エン (HAN{\middel}EN)\\\hline
円安 & yen + ease = \textit{weak yen} & エン{\middel}やす (EN{\middel}yasu)\\\hline
円周 & circle + around = \textit{circumference} & エン{\middel}シュウ (EN{\middel}SH\={U})\\\hline
\end{somme}

\begin{sentence}
\underline{あの} \ruby{\underline{子}}{こ}\ruby{は}{わ}\ruby{\underline{百円}}{ヒャクエン}\ruby{を}{お}\underline{なくした}。\\\hline
ano ko wa HYAKU{\middel}EN o nakushita.\\\hline
(that) (child) (one hundred yen) (lost)\\\hline
=\textit{That child lost a hundred yen.}\\\hline
\end{sentence}
\end{multicols}
\vspace*{-\baselineskip}
\vhrulefill{5pt}
\vspace*{\fill}
%%--------------------------------------------------------------------
%%--------------------------------------------------------------------
\pagebreak
\vspace*{\fill}
\begin{multicols}{2}
\begin{glyph}{Basis (元)}{basis}\label{glyph:basis}
This character expresses the idea of a basis or the original state of something.\\\hline
When used with its kun-yomi before words such as ``president'', it can be translated as ``ex'' or ``former''.
\end{glyph}

\begin{comps}
\comprtwocone & 二 + 儿 \\\hline
    \comprthreecone & Two (\ref{glyph:two}) + Man/A person (at the bottom of a
    character)\\\hline
\end{comps}

\begin{remglyph}
    \remcom{Two} \remcom{people} are the \hooglig{basis} of any
    relationship.\\\hline
\end{remglyph}

\begin{prons}
\pronsrtwocone & \textbf{ゲン、ガン (GEN, GAN)} & \textbf{もと (moto)}\\\hline
\pronsrthreecone & --- & --- \\\hline
\rye{3}{ガン only occurs in the first position with a few common words.}
\end{prons}

\begin{remprons}
\comon{Genghis Khan} holding the \comkun{motorists} at \comon{gunpoint} had no \hooglig{basis}.\\\hline
\end{remprons}

\begin{somme}
元々 & basis + basis = \textit{from the first} & もと{\middel}もと (moto{\middel}moto)\\\hline
元日 & basis + sun (day) = \textit{New Year's Day} & ガン{\middel}ジツ (GAN{\middel}JITSU)\\\hline
手元 & hand + basis = \textit{at hand} & て{\middel}もと (te{\middel}moto)\\\hline
元気 & basis + spirit = \textit{high spirits} & ゲン{\middel}キ (GEN{\middel}KI)\\\hline
火元 & fire + basis = \textit{origin of a fire} & ひ{\middel}もと (hi{\middel}moto)\\\hline
元来 & basis + come = \textit{by nature} & ガン{\middel}ライ (GAN{\middel}RAI)\\\hline
\end{somme}

\begin{sentence}
\ruby{\underline{元日}}{ガンジツ}な\underline{の}で\ruby{\underline{人}}{ひと}が\ruby{\underline{少}}{すく}\underline{ない}。\\\hline
GAN{\middel}JITSU na no de hito ga suku{\middel}nai.\\\hline
(New Year's Day) (because) (person) (few).\\\hline
=\textit{It's New Year's Day, so there aren't many people.}\\\hline
\end{sentence}
\end{multicols}
\vspace*{-\baselineskip}
\vhrulefill{5pt}
%%--------------------------------------------------------------------
%%--------------------------------------------------------------------
\begin{multicols}{2}
\begin{glyph}{Neck (首)}{neck}\label{glyph:neck}
Neck in the physical sense.\\\hline
As well as the ideas of chief/leader.
\end{glyph}

\begin{comps}
\comprtwocone & 䒑 + 自 \\\hline
\comprthreecone & grass + Self (\ref{glyph:neck})\\\hline
\end{comps}

\begin{remglyph}
I want to buy a \remcom{grass} \hooglig{necklace} \remcom{myself}.\\\hline
\end{remglyph}

\begin{prons}
\pronsrtwocone & \textbf{シュ (SHU)} & \textbf{くび (kubi)}\\\hline
\pronsrthreecone & --- & --- \\\hline
\end{prons}

\begin{remprons}
\comon{Shooting} \comkun{cubic} \hooglig{necks}.\\\hline
\end{remprons}

\begin{somme}
首 & \textit{neck} & くび (kubi)\\\hline
自首 & self + neck = \textit{surrender} & ジ{\middel}シュ (JI{\middel}SHU)\\\hline
手首 & hand + neck = \textit{wrist} & て{\middel}くび (te{\middel}kubi)\\\hline
足首 & leg + neck = \textit{ankle} & あし{\middel}くび (ashi{\middel}kubi)\\\hline
\end{somme}

\begin{sentence}
\ruby{\underline{高木}}{たかぎ}\underline{さん}\ruby{は}{わ}\ruby{\underline{首}}{くび}に\ruby{\underline{水}}{みず}\ruby{を}{お}\underline{かけました}。\\\hline
Takagi-san wa kubi ni mizu o kakemashita\\\hline
(Takagi-san) (neck) (water) (put)\\\hline
=\textit{Takagi-san put some water on his neck.}\\\hline
\end{sentence}
\end{multicols}
\vspace*{-\baselineskip}
\vhrulefill{5pt}
\vspace*{\fill}
%%--------------------------------------------------------------------
%%--------------------------------------------------------------------
\pagebreak
\vspace*{\fill}
\begin{multicols}{2}
\begin{glyph}{Walk (歩)}{walk}\label{glyph:walk}
Walk; Step; Pace.\\\hline
A secondary meaning of ``rate'' (as used in the financial sense) is found only with ブ, a less common on-yomi.
\end{glyph}

\begin{comps}
\comprtwocone & 止 + 少 \\\hline
\comprthreecone & Stop (\ref{glyph:stop}) + Few (\ref{glyph:small})\\\hline
\end{comps}

\begin{remglyph}
There are \remcom{few} \remcom{stops} in \hooglig{walking}.\\\hline
\end{remglyph}

\begin{prons}
\pronsrtwocone & \textbf{ホ (HO)} & \textbf{ある (aru)}\\\hline
\pronsrthreecone & ブ、フ (BU, FU) & あゆ (ayu) \\\hline
\end{prons}

\begin{remprons}
I \hooglig{walked} to the \comon{hospital} to receive \ucomkun{ayurvedic} treatment since I was \ucomon{bullied} to eat \comkun{arugula} which turned out to be a \ucomon{foofy}.\\\hline
\end{remprons}

\begin{somme}
歩く & \textit{to walk} & ある{\middel}く (aru{\middel}ku)\\\hline
一歩 & one + walk = \textit{one step} & イッ{\middel}ポ ()\\\hline
歩き回る & walk + rotate = \textit{to ramble about} & ある{\middel}き　まわ{\middel}る (aru{\middel}ki mawa{\middel}ru)\\\hline
歩道 & walk + road = \textit{sidewalk} & ホ{\middel}ドウ (HO{\middel}D\={O})\\\hline
歩行 & walk + go = \textit{walking} & ホ{\middel}コウ (HO{\middel}K\={O})\\\hline
歩行者 & walk + go + individual = \textit{pedestrian} & ホ{\middel}コウ{\middel}シャ (HO{\middel}K\={O}{\middel}SHA)\\\hline
\end{somme}

\begin{sentence}
\ruby{\underline{家}}{いえ}{\;}\underline{から}{\;}\ruby{\underline{山}}{やま}{\;}\underline{まで} \ruby{\underline{歩}}{ある}\underline{きました}。\\\hline
ie kara yama made aru{\middel}kimashita.\\\hline
(house) (from) (mountain) (until) (walked)\\\hline
=\textit{(We) walked from the house to the mountain.}\\\hline
\end{sentence}
\end{multicols}
\vspace*{-\baselineskip}
\vhrulefill{5pt}
%%--------------------------------------------------------------------
%%--------------------------------------------------------------------
\begin{multicols}{2}
\begin{glyph}{Child (子)}{child}\label{glyph:child}
Child.\\\hline
This character can also be used to suggest something small.
\end{glyph}

\begin{comps}
\comprtwocone & 子 \\\hline
\comprthreecone & Child (\ref{glyph:child})\\\hline
\end{comps}

\begin{remglyph}
---\\\hline
\end{remglyph}

\begin{prons}
\pronsrtwocone & \textbf{シ (SHI)} & \textbf{こ (ko)}\\\hline
\pronsrthreecone & ス (SU) & --- \\\hline
\rye{3}{Both of these readings can at times become voiced in the second position.\\\hline
Japanese parents often choose this kanji for a suffix when naming their children.  It is read シ in the case of males, and こ when applied to females.}
\end{prons}

\begin{remprons}
It \ucomon{suited} the \hooglig{child} to be \comon{shielded} on the \comkun{cobbled} road.\\\hline
\end{remprons}

\begin{somme}
子 & \textit{child} & こ (ko)\\\hline
女子 & woman + child = \textit{women's} & ジョ{\middel}シ (JO{\middel}SHI)\\\hline
男子 & man + child = \textit{men's} & ダン{\middel}シ (DAN{\middel}SHI)\\\hline
王子 & king + child = \textit{prince} & オウ{\middel}ジ (\={O}{\middel}JI)\\\hline
子牛 & child + cow = \textit{calf} & こ{\middel}うし (ko{\middel}ushi)\\\hline
電子 & electric + child = \textit{electron} & デン{\middel}シ (DEN{\middel}SHI)\\\hline
\end{somme}

\begin{sentence}
\ruby{\underline{五月}}{ゴガツ}に\ruby{\underline{女王}}{ジョオウ}と\ruby{\underline{王子}}{オウジ}が\ruby{\underline{米国}}{ベイコク}\ruby{へ}{え}\ruby{\underline{行}}{い}\underline{きました}。\\\hline
GO{\middel}GATSU ni JO{\middel}\={O} to \={O}{\middel}JI ga BEI{\middel}KOKU e i{\middel}kimashita.\\\hline
(May) (queen) (prince) (America) (went)\\\hline
=\textit{In May, the queen and prince went to America.}\\\hline
\end{sentence}
\end{multicols}
\vspace*{-\baselineskip}
\vhrulefill{5pt}
\vspace*{\fill}
%%--------------------------------------------------------------------
%%--------------------------------------------------------------------
\pagebreak
\vspace*{\fill}
\begin{multicols}{2}
\begin{glyph}{Like (好)}{like}\label{glyph:like}
Like; Good; Favorable.\\\hline
Nothing but positive feelings are associated with this character.
\end{glyph}

\begin{comps}
\comprtwocone & 女 + 子 \\\hline
    \comprthreecone & Woman (\ref{glyph:woman}) +  Child (\ref{glyph:child})\\\hline
\end{comps}

\begin{remglyph}
Woman likes child.\\\hline
\end{remglyph}

\begin{prons}
\pronsrtwocone & \textbf{コウ (K\={O})} & \textbf{す (su)}\\\hline
\pronsrthreecone & --- & この (kono) \\\hline
\end{prons}

\begin{remprons}
I \hooglig{like} \comkun{sudoku} because it \ucomkun{connotes} \comon{causality}.\\\hline
\end{remprons}

\begin{somme}
好き & \textit{like} & す{\middel}き (su{\middel}ki)\\\hline
大好き & large + like = \textit{to really like} & ダイ　す{\middel}き (DAI su{\middel}ki)\\\hline
人好き & person + like = \textit{amiability} & ひと　ず{\middel}き (hito zu{\middel}ki)\\\hline
話し好き & speak + like = \textit{talkative} & はな{\middel}し　ず{\middel}き (hana{\middel}shi zu{\middel}ki)\\\hline
友好 & friend + like = \textit{friendship} & ユウ{\middel}コウ (Y\={U}{\middel}K\={O})\\\hline
同好 & same + like = \textit{same tastes} & ドウ{\middel}コウ (D\={O}{\middel}K\={O})\\\hline
\end{somme}

\begin{sentence}
\ruby{\underline{山口}}{やまぐち}\underline{さん}\ruby{は}{わ}\underline{あそこ}の\ruby{\underline{女}}{おんな}の\ruby{\underline{人}}{ひと}が\ruby{\underline{大好}}{ダイす}\underline{き}です。\\\hline
Yama{\middel}guchi-san wa asoko no onna no hito ga DAI su{\middel}ki desu.\\\hline
(Yamaguchi-san) (over there) (woman) (person) (really likes)\\\hline
=\textit{Yamaguchi-san really likes that woman over there.}\\\hline
\end{sentence}
\end{multicols}
\vspace*{-\baselineskip}
\vhrulefill{5pt}
%%--------------------------------------------------------------------
%%--------------------------------------------------------------------
\begin{multicols}{2}
\begin{glyph}{Old (古)}{old}\label{glyph:old}
Old; Ancient; Antique.
\end{glyph}

\begin{comps}
\comprtwocone & 十 + 口 \\\hline
\comprthreecone & Ten (\ref{glyph:ten}) + Mouth (\ref{glyph:mouth})\\\hline
\end{comps}

\begin{remglyph}
Both the \remcom{scarecrow} and \remcom{vampire} are \hooglig{old} fictional characters.\\\hline
\end{remglyph}

\begin{prons}
\pronsrtwocone & \textbf{コ (KO)} & \textbf{ふる (furu)}\\\hline
\pronsrthreecone & --- & --- \\\hline
\end{prons}

\begin{remprons}
The \hooglig{old} \comkun{furuncular} \comon{cop}.\\\hline
\end{remprons}

\begin{somme}
古い & \textit{old} & ふる{\middel}い (furu{\middel}i)\\\hline
中古 & middle + old = \textit{secondhand} & チュウ{\middel}コ (CH\={U}{\middel}KO)\\\hline
最古 & most + old = \textit{oldest} & サイ{\middel}コ (SAI{\middel}KO)\\\hline
古語 & old + words = \textit{archaic word} & コ{\middel}ゴ (KO{\middel}GO)\\\hline
古来 & old + come = \textit{time honored} & コ{\middel}ライ (KO{\middel}RAI)\\\hline
古里 & old + hamlet = \textit{old hamlet} & ふる{\middel}さと (furu{\middel}sato)\\\hline
\end{somme}

\begin{sentence}
\ruby{\underline{山下}}{やました}\underline{さん}\ruby{は}{さ}\ruby{\underline{中古}}{チュウコ}の\ruby{\underline{品}}{しな}\ruby{を}{お} \underline{よく} \ruby{\underline{買}}{か}\underline{います}。\\\hline
Yama{\middel}shita-san wa CH\={U}{\middel}KO no shina o yoku ka{\middel}imasu.\\\hline
(Yamashita-san) (secondhand) (goods) (often) (buys)\\\hline
=\textit{Yamashita-san often buys secondhand goods.}\\\hline
\end{sentence}
\end{multicols}
\vspace*{-\baselineskip}
\vhrulefill{5pt}
\vspace*{\fill}
%%--------------------------------------------------------------------
%%--------------------------------------------------------------------
\pagebreak
\vspace*{\fill}
\begin{multicols}{2}
\begin{glyph}{Fire (火)}{fire}\label{glyph:fire}
Fire.
\end{glyph}

\begin{comps}
\comprtwocone & 火 \\\hline
\comprthreecone & Fire (\ref{glyph:fire})\\\hline
\end{comps}

\begin{remglyph}
---\\\hline
\end{remglyph}

\begin{prons}
\pronsrtwocone & \textbf{カ (KA)} & \textbf{ひ (hi)}\\\hline
\pronsrthreecone & --- & ほ (ho) \\\hline
\end{prons}

\begin{remprons}
A \ucomkun{hobby} of \comon{cussing} \comkun{heaps} of \hooglig{fire}.\\\hline
\end{remprons}

\begin{somme}
火 & \textit{fire} & ひ (hi)\\\hline
火元 & fire + basis = \textit{the origin of a fire} & ひ{\middel}もと (hi{\middel}moto)\\\hline
大火 & large + fire = \textit{conflagration} & タイ{\middel}カ (TAI{\middel}KA)\\\hline
火山 & fire + mountain = \textit{volcano} & カ{\middel}ザン (KA{\middel}SEI)\\\hline
花火 & flower + fire = \textit{fireworks} & はな{\middel}び (hana{\middel}bi)\\\hline
火曜日 & fire + day of the week + sun (day) = \textit{Tuesday} & カ{\middel}ヨウ{\middel}び (KA{\middel}Y\={O}{\middel}bi)\\\hline
火星 & fire + star = \textit{Mars} & カ{\middel}セイ (KA{\middel}SEI)\\\hline
\end{somme}

\begin{sentence}
\ruby{\underline{日本}}{ニホン}に\ruby{は}{わ}\ruby{\underline{火山}}{カザン}が\underline{あります}。\\\hline
NI{\middel}HON ni wa KA{\middel}ZAN ga arimasu.\\\hline
(Japan) (volcano) (are)\\\hline
=\textit{There are volcanoes in Japan.}\\\hline
\end{sentence}
\end{multicols}
\vspace*{-\baselineskip}
\vhrulefill{5pt}
%%--------------------------------------------------------------------
%%--------------------------------------------------------------------
\begin{multicols}{2}
\begin{glyph}{Winter (冬)}{winter}\label{glyph:winter}
Winter.
\end{glyph}

\begin{comps}
\comprtwocone & 夂 + 丶 + 丶 \\\hline
\comprthreecone & Overtaking + Dot + Dot\\\hline
\end{comps}

\begin{remglyph}
    \hooglig{Winter} \remcom{overtaking} two \remcom{dots}.\\\hline
\end{remglyph}

\begin{prons}
\pronsrtwocone & \textbf{トウ (T\={O})} & \textbf{ふゆ (fuyu)}\\\hline
\pronsrthreecone & --- & --- \\\hline
\end{prons}

\begin{remprons}
\comkun{Few usage} during a \comon{torturous} \hooglig{winter}.\\\hline
\end{remprons}

\begin{somme}
冬 & \textit{winter} & ふゆ (fuyu)\\\hline
冬休み & winter + rest = \textit{winter vacation} & ふゆ　やす{\middel}み (fuyu yasu{\middel}mi)\\\hline
立冬 & stand + winter = \textit{first day of winter} & リッ{\middel}トウ (RIT{\middel}T\={O})\\\hline
冬物 & winter + thing = \textit{winter clothing} & ふゆ{\middel}もの (fuyu{\middel}mono)\\\hline
冬空 & winter + empty = \textit{winter sky} & ふゆ{\middel}ぞら (fuyu{\middel}zora)\\\hline
\end{somme}

\begin{sentence}
\ruby{\underline{冬}}{ふゆ}に\underline{なる}と\underline{スイス}の\ruby{\underline{山}}{やま}\ruby{は}{わ}\ruby{\underline{美}}{うつく}\underline{しい}。\\\hline
fuyu ni naru to SUISU no yama wa utsuku{\middel}shii.\\\hline
(winter) (becomes) (Switzerland) (mountain) (beautiful)\\\hline
=\textit{The mountains of Switzerland are beautiful in winter.}\\\hline
\end{sentence}
\end{multicols}
\vspace*{-\baselineskip}
\vhrulefill{5pt}
\vspace*{\fill}
%%--------------------------------------------------------------------
%%--------------------------------------------------------------------
\pagebreak
